\paragraph{Adaptive digital baseline.}
On the previously evaluated common 500-source holdout, with three noise realizations per source, the later-added adaptive-downsampling BPG+LDPC baseline achieves higher PSNR than the proposed method at all six SNRs, by 2.665--3.575~dB. Conversely, the proposed method has lower LPIPS by 0.329--0.498, higher DINOv2-L cosine similarity by 0.237--0.410, and higher ConvNeXt source-prediction agreement by 39.6--55.2 percentage points. All corresponding pointwise 95\% source-paired confidence intervals exclude zero. At 13~dB, for example, the adaptive-baseline-minus-proposed PSNR difference is $+2.843$~dB, with a 95\% interval of $[2.759,2.930]$~dB. These results demonstrate a distortion--perception trade-off and do not support a claim of universal superiority for the proposed method.

The adaptive baseline also improves PSNR over continuous latent-space JSCC at every SNR, while its LPIPS and prediction agreement are worse throughout; their DINOv2-L difference at 19~dB remains inconclusive. Relative to adapted SwinJSCC-80k, adaptive BPG is worse on all four metrics at 1, 4, and 7~dB; its PSNR difference at 10~dB is inconclusive. At 13~dB it improves PSNR but worsens LPIPS, with DINOv2-L and agreement intervals including zero. It improves all four metrics at 19~dB, which is outside Swin's training and calibration range. Compared with native256 BPG+LDPC, adaptive BPG improves all four metrics at every SNR. No source overflow occurs; 15 of 1,500 frames fail the body CRC at 7~dB, and all other tested frames decode successfully. All failures remain included in the quality statistics.
